% Standalone independent audit; compile directly with pdfLaTeX twice.
% MAIN VERDICT: VERIFIED (mathematical statements).
% Exact P04 v1.0 and every supplied project file are preserved unchanged.
% Historical numerical-run statistics have a separately identified provenance limit.
\documentclass[11pt]{article}
\usepackage[T1]{fontenc}
\usepackage[utf8]{inputenc}
\usepackage{lmodern,amsmath,amssymb,amsthm}
\usepackage[margin=0.85in]{geometry}
\usepackage{array,booktabs,longtable,enumitem,needspace}
\usepackage[hidelinks]{hyperref}
\newcolumntype{L}[1]{>{\raggedright\arraybackslash}p{#1}}
\setlength{\emergencystretch}{3em}
\setlength{\parindent}{0pt}
\setlength{\parskip}{5pt}
\setlist{nosep,leftmargin=*}
\newcommand{\RR}{\mathbb R}
\newcommand{\CC}{\mathbb C}
\newcommand{\ZZ}{\mathbb Z}
\newcommand{\one}{\mathbf 1}
\newcommand{\Dom}{(-1,\infty)^n}
\newcommand{\diag}{\operatorname{diag}}
\newcommand{\sa}{\operatorname{sa}}
\newcommand{\lin}[1]{\textbf{L#1}}
\newcommand{\sha}[2]{\par{\small\ttfamily #1\allowbreak #2}\par}
\newcommand{\pass}{\textbf{PASS}}
\title{A04: Independent mathematical audit of P04\\Arbitrary-$(n,q)$ Fourier spectrum and normality}
\author{Independent reconstruction and adversarial review}
\date{27 September 2026}
\begin{document}
\maketitle

\section{Verdict, exact version, and audit boundary}
\textbf{MAIN VERDICT: VERIFIED.}
All substantive mathematical results in the exact P04 version identified
below pass independent reconstruction under their stated hypotheses. No
mathematical or meaning-changing typographical correction is required.
The optional propagator-norm and nonlinear-local-stability statements also
pass; no separate extended mathematical verdict is needed.

The proof being audited is
\path{proofs/P04_arbitrary_nq_fourier_spectrum.tex},
\textbf{P04 v1.0, dated 2026-09-27, DRAFT}, exactly
\textbf{25,792 bytes and 573 lines}. Its full-file SHA-256 is
\sha{ab80f695b3e20c852c572b7b174ecaf99}{64eac1aba62e3e55b2b4d3e60754e0c}
and its body SHA-256, from the first line-anchored section command through
EOF, including the final newline, is
\sha{3e3d58230c87af3b5f07069d61065362d}{c081d172ad3b236057fefda25fa5653}
The supplied \path{cyclic_network_project_P04_draft.zip} has SHA-256
\sha{11a2383d3f61e001c400f61744d6cacd}{94b0b02d670bd7240e0dbb948fc0610f}
Every P04 line reference below refers to that exact 573-line source,
including comments and blank lines.

The supplied archive, its latest manifest record, the current proof headers,
and the P01--P03 audit-closure records select the current versions.
In particular, the final P04 delivery hash above is authoritative; the
manifest's earlier pre-P04 input hash beginning \texttt{b68f6b7b} identifies
the former placeholder, not this audit target. All fifteen requested current
files were located and consulted before mathematical reconstruction.
The additional shared macro file was checked for notation.

\textbf{Evidence limitation U01 (historical computations only).}
The separately named \path{P04_validation_support.zip} is absent from the
supplied 75-member archive and was not supplied separately with this task.
Consequently the exact historical run counts and floating-point error
maxima in P04 \lin{530--545} have not been independently reproduced.
Those provenance assertions are marked \textbf{UNRESOLVED} in the line
review. They are not mathematical premises, and this audit does not
certify them by copying the manifest. The mathematical verdict rests on
the proofs below; new exact checks are separately identified. No theorem
or downstream mathematical dependency is conditional on the missing run.

\textbf{Preservation and closure.}
No P04 edit, status promotion, registry edit, compendium edit, or later-module
work is performed. Every original archive member is preserved byte-for-byte.
This audit verifies the specified version while its original DRAFT header
remains intact. \textbf{P05 is mathematically cleared after administrative
A04 closure records this exact-version verdict.} P05 is not begun here.

\section{Authenticated inputs and dependency boundary}
All hashes below cover complete file bytes, not normalized text. Paths are
relative to the supplied archive's \path{cyclic_network_project/} root.
The support file \path{A04_support/input_hashes.json} also records byte and
line counts. Historical status prose is read through the later closure
records, not mistaken for the current status.

\begin{longtable}{L{0.44\textwidth}L{0.50\textwidth}}
\toprule Current input & Full-file SHA-256\\\midrule\endhead
\small\path{project_manifest.md} & \small\texttt{58c5b57bd02f1e820af15db59296fce9}\newline\texttt{cd23887b33a1907f88f8e4727d406762}\\
\small\path{claim_registry.md} & \small\texttt{d546f1128362d7894cc9948174779490}\newline\texttt{95e65527b92289707188314ca935031b}\\
\small\path{frozen_theorem_package.md} & \small\texttt{724adf3d04a7f9661a2e53fbae359192}\newline\texttt{e8a5a5ee56737bdca2379209b556ea63}\\
\small\path{final_novelty_verdict.md} & \small\texttt{5fa8ddd5ada9628b104dba68dd535356}\newline\texttt{442c5e62ec45b72f3c2c354dec6037ea}\\
\small\path{02_theorem_dependency_map.tex} & \small\texttt{a1a4435f36fc1d1b6f0b2c5226a91328}\newline\texttt{3c76658d55b96493e6c50a8693831cac}\\
\small\path{00_canonical_model.tex} & \small\texttt{c209981cb2924e8a782f24aa2904e8c8}\newline\texttt{5a105ff99319dd6aada12f48a351c2a4}\\
\small\path{01_notation_and_assumptions.tex} & \small\texttt{97d993c6d3986524dbcc29659ee73d25}\newline\texttt{7be1b03a35c44401e8ef4534303eeb91}\\
\small\path{proofs/P01_wellposedness_and_absorbing_box.tex} & \small\texttt{0d6950821aa83f3d09e8bfe84dab2b0b}\newline\texttt{5020aa6e7881cf6928e984bff1890f1e}\\
\small\path{audits/A01_wellposedness_and_absorbing_box.tex} & \small\texttt{8d5c1b2ec782dbe49b07d2bbe529ee47}\newline\texttt{4d4c142695c53b523ccb41d63bb33098}\\
\small\path{proofs/P02_synchronous_subspace_and_equilibria.tex} & \small\texttt{8e057c2cd009659eee45906844e25f4f}\newline\texttt{38c156e5eb2e4130e824aa5e8299c51a}\\
\small\path{audits/A02_synchronous_subspace_and_equilibria.tex} & \small\texttt{ca3d723d9fcfac14f5129583441e611a}\newline\texttt{65e41848ec44f3b5e3a0072f795d84ac}\\
\small\path{proofs/P03_origin_and_linearization.tex} & \small\texttt{6c3c9bfb2307bdc043eb7f6a874cb222}\newline\texttt{e2ac2ce0750b339181f0088dd0bd859e}\\
\small\path{audits/A03_origin_and_linearization.tex} & \small\texttt{e1008f98ba9667b92b8d74f1bc7d21e0}\newline\texttt{21414b454ac976ae0db1500b4b7ef8d1}\\
\small\path{proofs/P04_arbitrary_nq_fourier_spectrum.tex} & \small\texttt{ab80f695b3e20c852c572b7b174ecaf9}\newline\texttt{964eac1aba62e3e55b2b4d3e60754e0c}\\
\small\path{proofs_compendium.tex} & \small\texttt{7acc9796e47c47a91d6e32b192dddc15}\newline\texttt{c8e0c164be3b8c59db201265ef8137b5}\\
\small\path{shared_macros.tex} & \small\texttt{f4a18a440e1214b619bf18f4a81fee50}\newline\texttt{7239460ca796d8ca3916791069d6a7ad}\\
\bottomrule
\end{longtable}

P01 v1.1's current hash agrees with its closure after the sole archived-radius
wording correction. P02 v1.1's current hash agrees with closure of its
ordered-root-limit wording correction. Its old comment that A02 source is
absent is superseded by the later manifest: the authenticated A02 source is
present. P03 v1.1's current hash agrees with its metadata-only A03 closure;
its body hash remains
\sha{07d76c8ce0b5232827e9f2c74245188b5}{5b0ce0d1732b3cb2e528d69e8d6b50a}
The A01/A02 original audit verdicts are not retroactively rewritten; their
required corrections have already been closed. No re-audit of P03's
differentiation is needed or undertaken.

The exact mathematical inputs are canonical indexing and positive parameters;
P03 \lin{268--325} for the homogeneous Jacobian and conditional equilibrium
simplification; P02's invariant synchronous restriction and nonzero equilibrium
identity; and P01 \lin{80--95} for smoothness and local ODE continuation.
Finite-dimensional geometric sums, linear algebra, differentiation,
variation of constants, and Gronwall complete the proof. P04 explicitly
invokes the standard $C^1$ linearized-instability theorem at \lin{444--451}.
Its statement and hypotheses are correct; Section~\ref{sec:nonlinear}
also proves the needed instance directly, so no unverified later module
or external literature claim is needed. In a strictly literal dependency
inventory that standard ODE theorem is an additional declared background
input, not a consequence of smoothness alone.

\section{Shift phases, including wraparound}
Let $n\ge5$, $1\le q<n$, $j,m\in\{0,\ldots,n-1\}$,
$\theta_m=2\pi m/n$, and
$v_j^{(m)}=n^{-1/2}e^{i\theta_m j}$. For any integer $\ell$, write
$j-\ell=r+kn$, where $0\le r<n$. Then
\[
 v^{(m)}_{(j-\ell)\bmod n}=n^{-1/2}e^{i\theta_m r}
 =n^{-1/2}e^{i\theta_m(j-\ell)}e^{-2\pi i mk}
 =e^{-i\ell\theta_m}v_j^{(m)}.
\]
Thus modular reduction produces no phase error, even for negative shifts.
In particular, at the predecessor wraparound $j=0$,
$v_{n-1}^{(m)}=e^{-i\theta_m}v_0^{(m)}$. At a successor-$q$
wraparound $j+q\ge n$, subtraction of $n$ produces the same factor one.
The three required phases are therefore
\[
 Pv^{(m)}=e^{-i\theta_m}v^{(m)},\qquad
 P^{-q}v^{(m)}=e^{iq\theta_m}v^{(m)},\qquad
 P^qv^{(m)}=e^{-iq\theta_m}v^{(m)}.
\]
This establishes the orientation gate independently of the proposed spectrum.
All three P04 phases pass. No downstream claim requires correction for
a phase reversal.

\section{DFT normalization, complete spectrum, and ordering}
Define $Q_{jm}=n^{-1/2}e^{2\pi imj/n}$. Direct calculation gives both products:
\begin{align*}
 (Q^*Q)_{\ell m}&=\frac1n\sum_{j=0}^{n-1}e^{2\pi i(m-\ell)j/n}
                      =\delta_{\ell m},\\
 (QQ^*)_{jk}&=\frac1n\sum_{m=0}^{n-1}e^{2\pi im(j-k)/n}
                      =\delta_{jk}.
\end{align*}
For unequal indices the ratio is a nontrivial $n$th root of unity, and the
finite geometric sum is zero; for equal indices there are $n$ unit terms.
This directly proves both unitarity identities, without importing a DFT
normalization from another convention. P04's shorter argument for $QQ^*=I$
from the square left inverse is also valid.

At every synchronous state $x=x^*\one$ with $x^*>-1$, put $s=1+x^*$.
Import the verified P03 result
\[
 J=dI+CP+GP^{-q}-HP^q,\quad
 d=a-2kx^*-\frac{cx^*}{s^2}-\frac{g}{s}-hx^*,\quad
 C=\frac c s,\quad G=\frac{gx^*}{s^2},\quad H=hx^*.
\]
Applying the independently established phases to each summand, rather than
comparing formulas, gives
\begin{equation}\label{eq:a04-spectrum}
 \lambda_m=d+Ce^{-i\theta_m}+Ge^{iq\theta_m}-He^{-iq\theta_m}.
\end{equation}
Since the columns of $Q$ are orthonormal, they constitute a basis of $\CC^n$.
The column equations $Jv^{(m)}=\lambda_mv^{(m)}$ give $JQ=Q\Lambda$, with
$\Lambda=\diag(\lambda_0,\ldots,\lambda_{n-1})$ in exactly the column order
$0,1,\ldots,n-1$. Hence
\[
 Q^*JQ=\Lambda,\qquad J=Q\Lambda Q^*,\qquad
 \det(zI-J)=\prod_{m=0}^{n-1}(z-\lambda_m).
\]
There is no conjugation of the displayed eigenvalues in $Q^*JQ$ and no
reversal of mode order. The negative exponent belongs to the forward
coordinate transform $Q^*$, while the basis matrix $Q$ has the positive
exponent. Every label supplies an eigenvector and all $n$ coordinates are
exhausted. Neither distinct offsets nor distinct eigenvalues is used.

\section{Dispersion, conjugacy, and singleton labels}
Euler's formulas in \eqref{eq:a04-spectrum} give
\begin{align*}
 r_m:=\Re\lambda_m&=d+C\cos\theta_m+(G-H)\cos(q\theta_m),\\
 \omega_m:=\Im\lambda_m&=-C\sin\theta_m+(G+H)\sin(q\theta_m).
\end{align*}
Specifically, $-He^{-iq\theta}=-H\cos(q\theta)+iH\sin(q\theta)$.
The critical $+H$ imaginary contribution is correct. The signed $\omega_m$
need not be positive and is distinct from the reserved future Hopf frequency
$\omega>0$.

Let $\bar m=(-m)\bmod n$ in the stated label set. Then
$\theta_{\bar m}\equiv-\theta_m\pmod{2\pi}$; all coefficients are real.
It follows componentwise that
\[
 v^{(\bar m)}=\overline{v^{(m)}},\quad
 \lambda_{\bar m}=\overline{\lambda_m},\quad
 r_{\bar m}=r_m,\quad \omega_{\bar m}=-\omega_m.
\]
At $m=0$ this uses label zero, not a second label $n$. The overline in
P04 \lin{130--131} covers the entire Fourier vector or eigenvalue; the bar
in the subscript is the previously defined modular label operation.

Self-conjugacy means $n\mid2m$. With $0\le m<n$, this gives just $m=0$
for odd $n$, and $m=0,n/2$ for even $n$. Thus
\[
 \lambda_0=d+C+G-H,\qquad
 \lambda_{n/2}=d-C+(G-H)(-1)^q\quad(n\text{ even}).
\]
Both are real. The second equals $d-C+G-H$ for even $q$ and
$d-C-G+H$ for odd $q$. When $q=n/2$, its factor is $(-1)^{n/2}$,
as obtained by substituting $m=n/2$ in the half-offset formula.
There are $(n-1)/2$ nonsingleton pairs for odd $n$ and $n/2-1$ for even $n$.
Conversely a nonsingleton pair is allowed to have $\omega_m=0$, yielding
one repeated real numerical eigenvalue. Conjugate labels, equal numerical
eigenvalues, and equal real parts are three different relations. In particular,
an equality of real parts does not impose equality of imaginary parts.

\section{Numerical multiplicity and semisimplicity}
For a numerical eigenvalue $\eta$, let
$M_\eta=\{m:\lambda_m=\eta\}$. Write $z=Qw$. Then
\[
 (J-\eta I)z=0\quad\Longleftrightarrow\quad
 (\lambda_m-\eta)w_m=0\quad\text{for all }m.
\]
Exactly $|M_\eta|$ coordinates are free. Therefore
\[
 \ker(J-\eta I)=\operatorname{span}_{\CC}\{v^{(m)}:m\in M_\eta\},
 \qquad \dim_{\CC}\ker(J-\eta I)=|M_\eta|.
\]
The characteristic product proves that the algebraic multiplicity is also
$|M_\eta|$. The diagonal form has no nontrivial Jordan block, so every
eigenvalue is semisimple. This proof is unchanged by a collision of
nonconjugate labels, a real value from a conjugate pair, or cancellation of
matrix coefficients. A singleton label is not an assertion of a simple
numerical eigenvalue.

The draft's example checks exactly: $q=1$, $x^*=1$, and
$(a,k,c,g,h)=(2,1,4,4,1)$ give $f(1)=2-2+(4-4)/2=0$,
$d=-4$, $C=2$, $G=H=1$. Thus $J=-4I+P+P^{-1}$ and
$\lambda_m=-4+2\cos\theta_m$. Each nonsingleton conjugate pair contributes
a real value with multiplicity two; cosine equality on this grid occurs
only at the same or conjugate label. The origin is different: $G=H=0$,
$C=c>0$, $d=a-g$, and the $n$ distinct roots of unity give simple values
$a-g+ce^{-i\theta_m}$, consistent with P03.

\section{Every offset collision}
All equalities below are exact identities of matrices and of all labeled
eigenvalues, not limits of distinct-offset formulas.
\begin{enumerate}[label=(\roman*)]
\item $q=1$: the $C$ and $-H$ contributions occupy the same shift:
\[
 J=dI+(C-H)P+GP^{-1},\qquad
 \lambda_m=d+(C-H)e^{-i\theta_m}+Ge^{i\theta_m}.
\]
Taking parts gives $r_m=d+(C+G-H)\cos\theta_m$ and
$\omega_m=(-C+G+H)\sin\theta_m$.
\item $q=n-1$: $P^{-q}=P$ and $P^q=P^{-1}$, so
\[
 J=dI+(C+G)P-HP^{-1},\qquad
 \lambda_m=d+(C+G)e^{-i\theta_m}-He^{i\theta_m}.
\]
Independently, $e^{iq\theta_m}=e^{-i\theta_m}$ and
$e^{-iq\theta_m}=e^{i\theta_m}$ reduce the master spectrum to the same
expression. Here $r_m=d+(C+G-H)\cos\theta_m$ and
$\omega_m=-(C+G+H)\sin\theta_m$.
\item Even $n$, $q=n/2$: $P^{-q}=P^q$ and
$e^{\pm iq\theta_m}=(-1)^m$. Therefore
\[
 J=dI+CP+(G-H)P^{n/2},\qquad
 \lambda_m=d+Ce^{-i\theta_m}+(G-H)(-1)^m.
\]
Its parts are $r_m=d+C\cos\theta_m+(G-H)(-1)^m$ and
$\omega_m=-C\sin\theta_m$.
\end{enumerate}
P03's exhaustive collision classification is used only to identify which
offsets meet. In all cases contributions add; none is overwritten.
If $C-H=0$, $C+G=0$, or $G-H=0$ where applicable, the corresponding
term simply vanishes. The basis, spectrum formula, and multiplicity
argument remain valid. These cases do not create a missing eigenvector.

\section{Normality, in two independent forms}
From unitary diagonalization,
\[
 JJ^*=Q\Lambda\Lambda^*Q^*
 =Q\diag(|\lambda_m|^2)Q^*
 =Q\Lambda^*\Lambda Q^*=J^*J.
\]
Alternatively, $P$ permutes an orthonormal coordinate basis, so
$P^*=P^{-1}$. The coefficients are real, and consequently
\[
 J^*=dI+CP^{-1}+GP^q-HP^{-q}.
\]
Every integer power of $P$ commutes with every other. Expanding $JJ^*$
and $J^*J$ termwise therefore proves equality. Negative powers could
equivalently be replaced by nonnegative powers modulo $n$.

Neither proof uses spectral simplicity, positive $G,H$, nonzero combined
coefficients, separate offsets, or coprimality. They cover negative
synchronous states $-1<x^*<0$, where $G,H<0$. The object remains a
homogeneous-state Jacobian. P04 \lin{266--267} explicitly withholds normality
at general nonsynchronous or heterogeneous states; nothing here promotes it.

\section{Synchronous tangent mode and equilibrium specialization}
Differentiating the verified scalar field directly gives
\[
 f(r)=ar-(k+h)r^2+\frac{(c-g)r}{1+r},\qquad
 f'(r)=a-2(k+h)r+\frac{c-g}{(1+r)^2}.
\]
The quotient derivative is $d[r/(1+r)]/dr=1/(1+r)^2$.
Using $s-x^*=1$, the unsimplified coefficients yield
\begin{align*}
 d+C+G-H
 &=a-2(k+h)x^*+
   c\left(\frac1s-\frac{x^*}{s^2}\right)
   +g\left(-\frac1s+\frac{x^*}{s^2}\right)\\
 &=a-2(k+h)x^*+\frac{c-g}{s^2}=f'(x^*).
\end{align*}
This uses no equilibrium relation. Independently, differentiating
$F(r\one)=f(r)\one$ gives $DF(r\one)\one=f'(r)\one$;
$v^{(0)}=\one/\sqrt n$ spans the synchronous tangent space.
Along a nonstationary synchronous trajectory the tangent coefficient is
$f'(r(t))$; no autonomous equilibrium stability claim is made at that state.

Only now let $x^*\ne0$ be an equilibrium and put $A=k+h$. P02 gives
$a-Ax^*+(c-g)/s=0$. Substitution into $d$ gives
\begin{align*}
 d&=Ax^*-\frac{c-g}{s}-2kx^*-\frac{cx^*}{s^2}
                          -\frac gs-hx^*\\
  &=-kx^*-\frac cs-\frac{cx^*}{s^2}
   =-kx^*-\frac{c(1+2x^*)}{s^2}.
\end{align*}
With $B(r)=a-Ar+(c-g)/(1+r)$, $f=rB$ and $B(x^*)=0$,
\begin{align*}
 \lambda_0&=x^*B'(x^*)=-x^*\left(A+\frac{c-g}{s^2}\right)\\
 &=-\frac{x^*}{s}\bigl(As+Ax^*-a\bigr)
   =\frac{x^*}{s}(a-A-2Ax^*).
\end{align*}
The second line uses $c-g=(Ax^*-a)s$. All formulas pass. They apply to a
positive equilibrium as requested and algebraically to a negative nonzero
equilibrium in $U$; only in the positive case are $G,H>0$.
They are not licensed by an arbitrary synchronous state. The origin is
retained separately with $\lambda_0=a+c-g$; one must not infer $B(0)=0$
merely from $f(0)=0$.

\section{Exact propagator norm and linear stability}
For the equilibrium linear system, the already complete spectrum gives
$\sa(J)=\max_m r_m$. The exponential series commutes with the finite
unitary change of basis, so
\[
 e^{tJ}=Q\diag(e^{t\lambda_m})Q^*.
\]
For a complex diagonal matrix $D=\diag(d_m)$,
$\|Dw\|_2^2\le(\max_m|d_m|)^2\|w\|_2^2$, with equality on a maximizing
coordinate vector. Since $Q$ and $Q^*$ preserve the Euclidean norm,
\begin{equation}\label{eq:a04-norm}
 \|e^{tJ}\|_{2,\CC}=\max_m|e^{t\lambda_m}|
 =\max_m e^{tr_m}=e^{t\sa(J)},\qquad t\ge0.
\end{equation}
The time restriction is correct: for negative $t$ the maximum would involve
the smallest real part. Repeated values cause no change in this argument;
a purely imaginary value contributes modulus one.

For every real matrix $T$, write a complex vector as $z=u+iv$ with real
$u,v$. If $K=\|T\|_{2,\RR}$, then
\[
 \|Tz\|_2^2=\|Tu\|_2^2+\|Tv\|_2^2
 \le K^2(\|u\|_2^2+\|v\|_2^2)=K^2\|z\|_2^2.
\]
Thus the complex norm is at most $K$. It is at least $K$ because real
vectors are among the complex vectors. Hence the norms coincide; apply
this to the real matrix $e^{tJ}$. P04's argument is correct, including
the otherwise implicit reverse inequality.

If $\sa(J)=-\delta<0$, \eqref{eq:a04-norm} gives an exponential estimate
with $M=1$ and decay exponent $\delta$. Conversely, if an exponential
estimate with $M\ge1$, $\beta>0$ holds for every initial state, the induced
norm satisfies $e^{t\sa(J)}\le Me^{-\beta t}$. Taking logarithms for $t>0$
and then $t\to\infty$ gives $\sa(J)\le-\beta<0$. This proves both
directions directly. The argument applies in both real and complex spaces.

\subsection{Real blocks, zero abscissa, and absence of Jordan growth}
For a nonsingleton pair, let $v=v^{(m)}$ and form
$a_m=\sqrt2\Re v$, $b_m=\sqrt2\Im v$. They are real orthonormal vectors:
$a_m=(v+\bar v)/\sqrt2$ and $b_m=(v-\bar v)/(i\sqrt2)$, and $v,\bar v$
are distinct orthonormal Fourier columns. From $Jv=(r_m+i\omega_m)v$,
\[
 Ja_m=r_ma_m-\omega_mb_m,\qquad
 Jb_m=\omega_ma_m+r_mb_m.
\]
Thus in the ordered real basis $(a_m,b_m)$ the matrix and exponential are
\[
 \begin{pmatrix}r_m&\omega_m\\-\omega_m&r_m\end{pmatrix},\qquad
 e^{tr_m}\begin{pmatrix}\cos(t\omega_m)&\sin(t\omega_m)\\
                  -\sin(t\omega_m)&\cos(t\omega_m)\end{pmatrix}.
\]
The sign in P04 is correct for its positive-exponent Fourier convention.
Each singleton has a real unit eigenvector. These real bases together are
an orthonormal basis of $\RR^n$.

If $\sa(J)=0$, the full propagator norm is one, proving Lyapunov stability.
A maximizing singleton with eigenvalue zero is a constant nondecaying
solution. A maximizing nonsingleton pair with $\omega_m\ne0$ rotates
without decay; with $\omega_m=0$ its two-dimensional block is zero and
both coordinates are constant. Additional repeated real zero values add
constant directions, never Jordan chains. Hence attraction fails in every
case. If $\sa(J)>0$, an expanding real block gives solutions whose norms
grow as $e^{tr_m}$ from arbitrarily small initial data, proving linear
Lyapunov instability. These facts also prove that linear asymptotic
stability here is equivalent to exponential asymptotic stability.

\section{Hyperbolicity is stronger than invertibility}
The complete diagonalization gives separately
\[
 \text{hyperbolic}\ \Longleftrightarrow\ r_m\ne0\ \text{for all }m,
 \qquad J\text{ invertible}\ \Longleftrightarrow\ \lambda_m\ne0\ \text{for all }m.
\]
The first condition implies the second. The reverse is false within the
canonical family itself. At the origin with $n=6$, $a=2$, $c=1$, $g=5/2$
and arbitrary positive $k,h$ and canonical $q$,
\[
 \lambda_m=-\tfrac12+e^{-i\pi m/3},\qquad
 \lambda_1=-i\tfrac{\sqrt3}{2},\quad
 \lambda_5=i\tfrac{\sqrt3}{2}.
\]
No value is zero because a unit-modulus complex number cannot be $1/2$.
The equilibrium is invertible and nonhyperbolic; its $m=0$ eigenvalue
$1/2$ also shows that nonhyperbolicity can coexist with instability.
P04 never identifies zero real part with a zero eigenvalue. Hyperbolicity
alone likewise does not imply stability.

\section{Nonlinear local stability and instability}\label{sec:nonlinear}
All assertions in this section are at a synchronous equilibrium interior
to $U$. P01 supplies a $C^\infty$ vector field there. With
$y=x-x^*\one$, differentiability and $F(x^*\one)=0$ give
\[
 F(x^*\one+y)=Jy+R(y),\qquad R(y)=o(\|y\|_2).
\]

If $\sa(J)=-\delta<0$, choose $\rho>0$ so that the closed ball
$\overline B_\rho(0)$ is contained in the translated domain and
$\|R(y)\|_2\le(\delta/2)\|y\|_2$ there. The little-$o$ estimate and
openness of $U$ permit both requirements. While a solution is in the ball,
variation of constants and \eqref{eq:a04-norm} give
\[
 \|y(t)\|_2\le e^{-\delta t}\|y(0)\|_2+
 \frac\delta2\int_0^t e^{-\delta(t-u)}\|y(u)\|_2\,du.
\]
Multiply by $e^{\delta t}$ and apply Gronwall to obtain
$\|y(t)\|_2\le e^{-\delta t/2}\|y(0)\|_2$.
If $\|y(0)\|_2<\rho$, a first exit at norm $\rho$ contradicts this
bound. The orbit lies in the smaller compact ball of radius
$\|y(0)\|_2$ inside $U$; P01 local continuation excludes a finite maximal
forward time. This proves local exponential asymptotic stability and
validates every step of P04's bootstrap argument.

For $\sa(J)>0$, the invoked finite-dimensional $C^1$ instability theorem
requires an open neighborhood, an equilibrium, differentiability, and at
least one eigenvalue with positive real part. It allows other eigenvalues
on the imaginary axis. All hypotheses are checked in P04.
For completeness, the following direct argument proves the needed instance
without citing that theorem.

Use the real orthogonal block basis above. Let $E_+$ contain all blocks
with positive real part and $E_{\le0}$ the remaining blocks; write
$y=u+w$ in this orthogonal splitting. Let
$\alpha=\min\{r_m:r_m>0\}>0$ and choose a closed ball compactly contained in the translated domain where
$\|R(y)\|_2\le\varepsilon\|y\|_2$, with
$\varepsilon\le\alpha/(4\sqrt2)$. Define
$V=\|u\|_2^2-\|w\|_2^2$. Inside the cone $\|w\|_2\le\|u\|_2$,
the block real parts and the orthogonal projections $R_+,R_{\le0}$ imply
\begin{align*}
 \dot V
 &=2\langle u,J_+u\rangle-2\langle w,J_{\le0}w\rangle
       +2\langle u,R_+\rangle-2\langle w,R_{\le0}\rangle\\
 &\ge 2\alpha\|u\|_2^2
       -2(\|u\|_2+\|w\|_2)\varepsilon\|y\|_2\\
 &\ge (2\alpha-4\sqrt2\varepsilon)\|u\|_2^2
 \ge\alpha\|u\|_2^2\ge\alpha V.
\end{align*}
On the nonzero cone boundary the inequality is strictly positive, so
a trajectory starting with $w(0)=0$, $u(0)\ne0$ cannot leave the cone
before leaving the ball. Uniqueness excludes arrival at the equilibrium
in finite time. If the solution stayed in that ball forever, continuation
would give all forward times and $V(t)\ge e^{\alpha t}V(0)$, contradicting
boundedness. It must therefore leave a fixed neighborhood. Initial nonzero
$u(0)$ can be arbitrarily small, proving Lyapunov instability. Empty
$E_{\le0}$ simply means $w=0$. Zero-real-part blocks are permitted in
$E_{\le0}$ throughout; hyperbolicity is unnecessary.

If $\sa(J)=0$, the linear conclusion does not determine nonlinear stability
in general. The smooth scalar fields $\dot y=-y^3$ and $\dot y=y^3$ have
the same normal Jacobian zero. In the first,
$y(t)=y(0)/\sqrt{1+2y(0)^2t}$ tends to zero and stays no farther from
zero than its initial value; in the second positive small data leave a
fixed small neighborhood. These are valid illustrations of the limit of
linearization, not examples of DM parameter classification.

\subsection{Origin and state-space boundary}
The origin is interior to $U=(-1,\infty)^n$ but lies on the boundary of
$\RR_{\ge0}^n$. An instability direction supplied by a theorem on $U$
need not lie in a specified restricted state set. P04 \lin{461--467}
keeps that distinction explicit. At a positive equilibrium a sufficiently
small ball lies in the strict positive orthant, so the same local conclusions
apply there. P04 does not silently transfer origin instability to the
nonnegative state space. No claim of global attraction, basin size,
nonlinear boundary classification, or necessity of a Hurwitz Jacobian
for nonlinear asymptotic stability is made.

\section{Step-\texorpdfstring{$q$}{q} cycles and coprimality}
Write $\kappa=\gcd(n,q)$, $n=\kappa n'$, $q=\kappa q'$, with
$\gcd(n',q')=1$. The orbit of $j$ under addition of $q$ returns first at
the least positive $L$ such that $n\mid Lq$. This is equivalent to
$n'\mid Lq'$, hence to $n'\mid L$, so $L=n'=n/\kappa$.
Addition of $q$ preserves the residue modulo $\kappa$. Each residue class
contains $n/\kappa$ nodes, equal to the orbit size just proved; the orbit
therefore fills that class. There are exactly $\kappa$ disjoint cycles.
Subtracting $q$ traverses the same cycles in reverse.

This description concerns the step-$q$ permutation alone. The canonical
field also has the step-one predecessor coupling with $C=c/s>0$.
For $\kappa>1$, the two $q$ shifts remain inside residue classes whereas
the step-one shift connects distinct classes; it cannot be cancelled by
those within-class shifts. Thus the claimed step-$q$ partition does not
give a full Jacobian block decomposition by residue class. The case
$\kappa=1$ has only one class anyway. Fourier completeness, conjugacy,
multiplicities, normality, and the all-mode stability criterion used no
coprimality condition. Actual eigenvalues and their ordering can still
depend on $n,q$; independence of a proof hypothesis is not independence
of spectral values.

\section{Exact special cases and independent corroboration}
The following named examples were constructed independently. They are
checks of P04 identities, not a selection theorem, a branch analysis, or
evidence replacing the universal proofs.

For a single positive-equilibrium control take
$x^*=1$, $(a,k,c,g,h)=(5/2,1,1,2,1)$. The equilibrium residual is
$5/2-2+(1-2)/2=0$, and
\[
 (d,C,G,H)=(-7/4,1/2,1/2,1),\qquad
 \lambda_m=-7/4+\tfrac12z^{-1}+\tfrac12z^q-z^{-q},
 \quad z=e^{2\pi im/n}.
\]
This yields the following exact root-of-unity expressions, without
decimal trigonometric evaluation:
\begin{center}
\begin{tabular}{lll}
\toprule $(n,q)$ & Exact labeled eigenvalue & Step-$q$ cycles\\
\midrule
$(5,2)$ & $-7/4+z^{-1}/2+z^2/2-z^{-2}$ & one of length 5\\
$(6,2)$ & $-7/4+z^{-1}/2+z^2/2-z^{-2}$ & two of length 3\\
$(6,3)$ & $-7/4+z^{-1}/2-(-1)^m/2$ & three of length 2\\
$(7,1)$ & $-7/4-z^{-1}/2+z/2$ & one of length 7\\
$(8,4)$ & $-7/4+z^{-1}/2-(-1)^m/2$ & four of length 2\\
$(n,n-1)$ & $-7/4+z^{-1}-z$ & one of length $n$\\
\bottomrule
\end{tabular}
\end{center}
Each expression follows both from the directly assembled modular matrix
and from the independently derived phases; the support checker verifies
the equality exactly for every label in its stated finite cases.

Additional stress cases are:
\begin{enumerate}[label=(\arabic*)]
\item \textbf{Negative nonequilibrium state.}
At $x^*=-1/2$ and $(a,k,c,g,h)=(2,1,4,4,1)$,
$(d,C,G,H)=(7/2,8,-8,-1/2)$. The state lies in $U$; no equilibrium-only
simplification is used. Real coefficients still give both normality and
the same DFT proof.
\item \textbf{Repeated real value.} The draft's $q=1$ example gives
$J=-4I+P+P^{-1}$ and multiplicity two on nonsingleton pairs, as proved above.
\item \textbf{Nonconjugate coincident labels.}
With $n=6,q=3,x^*=1$ and $(a,k,c,g,h)=(4,1,4,8,1)$,
the equilibrium residual is zero and $(d,C,G,H)=(-4,2,2,1)$.
The labels $1,2$ both give $-4-i\sqrt3$, while $4,5$ both give
$-4+i\sqrt3$; the remaining values are $\lambda_0=-1$ and $\lambda_3=-7$.
The two nonreal values each have algebraic and geometric multiplicity two.
\item \textbf{First-offset cancellation.}
At $q=1,x^*=1,(a,k,c,g,h)=(3,1,2,4,1)$, the equilibrium residual is zero,
$C=H=1$, $G=1$, $d=-5/2$, and $J=-\tfrac52I+P^{-1}$.
\item \textbf{Half-offset cancellation.}
At even $n,q=n/2,x^*=1,(a,k,c,g,h)=(7/2,1,1,4,1)$,
$G=H=1$, $C=1/2$, $d=-7/4$, and $J=-\tfrac74I+\tfrac12P$;
the state is an equilibrium.
\item \textbf{Last-offset cancellation at a negative state.}
At $q=n-1$, $x^*=-1/2$, and all five parameters equal to one,
$(d,C,G,H)=(5/2,2,-2,-1/2)$, hence
$J=\tfrac52I+\tfrac12P^{-1}$. This is a state check, not an equilibrium claim.
\item \textbf{Zero abscissa with a nonzero imaginary pair.}
At $n=6,q=2,x^*=1,(a,k,c,g,h)=(17/4,1,1,1,13/4)$,
$f(1)=0$ and $(d,C,G,H)=(-7/4,1/2,1/4,13/4)$.
The real parts are $r_0=-17/4$, $r_1=r_5=0$,
$r_2=r_4=-1/2$, $r_3=-21/4$, and
$\lambda_1=i3\sqrt3/2$, $\lambda_5=-i3\sqrt3/2$.
This is an invertible, linearly stable but nonattracting equilibrium
linearization with norm identically one. No crossing or Hopf assertion follows.
\item \textbf{Zero abscissa with a real zero.}
At $x^*=1/4,(a,k,c,g,h)=(3,1,1,33/8,1)$, the equilibrium relation holds and
$(d,C,G,H)=(-121/100,4/5,33/50,1/4)$.
Here $\lambda_0=0$. Since $C>0$ and $G-H=41/100>0$,
every $m\ne0$ has $r_m<d+C+G-H=0$ from $\cos\theta_m<1$ and
$\cos(q\theta_m)\le1$. This checks a constant singleton critical direction;
it does not classify a bifurcation.
\end{enumerate}

\subsection{Retained exact-check execution}
The independently written standard-library checker in
\path{A04_support/checks/independent_exact_p04.py} completed successfully.
Its retained report records \textbf{39} admissible $(n,q)$ pairs,
\textbf{285} rational matrices, \textbf{2,242} eigenvalue occurrences,
\textbf{62,094} named exact assertions, \textbf{72} direct eigenspace-nullity
checks, and \textbf{four} formal multivariate polynomial identities.
Every named example above is included, with the general last-offset example
represented at $(n,q)=(6,5)$; the main families also cover every last offset
in dimensions five through ten. Positive and negative nonzero equilibria,
nonstationary states of both signs, and the origin are represented.

All arithmetic uses Python \texttt{Fraction}. A custom polynomial quotient
$\mathbb Q[z]/(\Phi_n(z))$ represents $z=e^{2\pi i/n}$ exactly; no numerical
eigensolver is used. Matrices are assembled by adding the four P03
contributions at modular node-column indices. Exact matrix-vector products
check every labeled eigenvector. Using $U=\sqrt n Q$ removes square roots
from both DFT orthogonality checks; inverse Fourier sums reconstruct every
matrix entry. Rational matrix multiplication checks normality. A direct
Faddeev--LeVerrier characteristic polynomial is compared coefficientwise
with the product of all Fourier factors. Selected exact Gaussian elimination
checks nullities, including real repetitions and nonconjugate coincidences.
Integer-shift checks include wraparounds from $-2n$ through $2n$.

The separate formal polynomial identities clear powers of $1+x^*$ and
check the all-state synchronous derivative and three equilibrium reductions
modulo the nonzero-equilibrium equation. They are identities, not sampled
substitutions. The finite matrix checks do not prove arbitrary-$n$ statements,
and neither exact time-dependent norms nor nonlinear behavior is inferred
from computation. Those assertions are proved analytically in this audit.
The checker reads P04 only to authenticate its SHA-256, never imports an
attached executable, and uses no old P04 internal computation as evidence.
Full exact inputs, characteristic polynomials, eigenvalue groups, cycle
partitions, and check counts are retained in its JSON report.

\section{Adversarial counterexample search}
The strongest quantifiers were challenged as follows. ``Proved'' means the
analytic argument above settles all inputs under the hypotheses, not just
that no sampled failure was observed.
\begin{longtable}{L{0.25\textwidth}L{0.66\textwidth}}
\toprule Claim challenged & Test and logical resolution\\\midrule\endhead
Multiplicity equals label count & Real pairs, nonconjugate coincidences, and
cancellations were tested. Free Fourier coordinates and the characteristic
product prove the equality and semisimplicity for every case.\\
Exact propagator norm & Repeated values and imaginary-axis blocks cannot
alter the diagonal induced norm. Unitarity proves equality for all $t\ge0$.
The claim would not follow from eigenvalues alone for a nonnormal matrix;
P04 supplies the needed normality.\\
Real/complex norm equality & The decomposition $z=u+iv$ gives the upper
bound and inclusion of real vectors gives the reverse bound, for every real
matrix. No eigenvector-realness assumption is needed.\\
Stability at zero abscissa & Imaginary rotations, zero singletons, and
repeated real zeros are bounded and nondecaying. Unitary diagonalization
excludes every Jordan growth mechanism. A nilpotent Jordan block would be
a counterexample outside the established normal class, not to P04.\\
Hyperbolicity criterion & The exact origin example is invertible with
nonzero purely imaginary eigenvalues. It refutes the stronger erroneous
equivalence with invertibility and confirms P04's distinction.\\
Collision reductions & All three matrix identities were checked before
taking phases. Vanishing combined coefficients do not change the diagonalizing
basis or remove labels.\\
Coprimality independence & $(6,2),(6,3),(8,4)$ retain a full DFT basis.
The geometric sum uses the entire grid, not the orbit length of the
step-$q$ subgraph. The proofs cover every allowed gcd.\\
Equilibrium restrictions & Negative nonequilibrium states were checked with
unsimplified coefficients; nonzero equilibrium identities were used only
with a vanishing scalar residual. The origin was checked separately.\\
Nonlinear instability with center spectrum & The cone proof uses all
positive blocks and allows zero real parts in the complementary subspace.
It proves the required conclusion without hyperbolicity.\\
\bottomrule
\end{longtable}
No mathematical counterexample was found. More strongly, the universally
quantified claims listed above have been analytically proved under the
actual hypotheses. Finite exact checks only corroborate those proofs.
Historical numerical-run details remain a different, unresolved provenance
question (U01), rather than a candidate mathematical counterexample.

\section{Terminology, scope, and downstream dependencies}
At a fixed parameter point the dominant/selected-label set is
$\operatorname*{arg\,max}_m r_m$, including all ties. Its elements may all
have negative growth rates. A first-instability statement additionally
requires a named equilibrium branch, an oriented parameter path, a prior
stable interval, and an actual sign change to positive spectral abscissa.
Along a continuous branch/path with a finite limiting equilibrium, the
finite family of continuous $r_m$ makes the maximum continuous; a boundary
of the strictly stable region therefore has maximum zero. Touching zero
alone does not establish an instability crossing. P04 includes these
conditions and assumes neither uniqueness nor transversality.

The exact source has no $q=2$ quadratic reduction in $\cos\theta$, no
discrete selection classification, no tie boundaries, no first-instability
regions, no large-$n$ selected wavenumber, and no five-node golden-ratio
formulas. Merely listing $(5,2)$ among finite checks does not enter P07.
It asserts no Hopf bifurcation, transversality, first Lyapunov coefficient,
or periodic branch. Its nonlinear statements are local stability results.
The optional norm identity supports, but does not perform, P17 response
analysis. The frozen supporting/known role and novelty verdict remain intact.

No proof refers to a theorem from P05 or later. All cross-references in
the mathematical argument lead to notation, P01--P03, or earlier P04
results. The standard ODE background was explicitly audited above.
The dependency graph is acyclic and its historical all-TODO prose is not
the current completion record.

\begin{itemize}
\item \textbf{P05:} after A04 closure, it may use the verified labeled
spectrum and dispersion, including both singleton types and all ties.
It must supply its own selection and branch/path crossing arguments.
\item \textbf{P08:} P04 clears the DM spectral input only. The comparator's
Jacobian and branch-specific comparison still need their own derivation;
P05 remains another dependency under the project map. No DM spectrum is
automatically a comparator spectrum.
\item \textbf{P11:} the signed frequency, label conjugacy, and numerical
multiplicity bookkeeping are available after closure. P11 must still
establish a simple numerical critical pair, a nonzero chosen positive
frequency, the remaining spectral inequalities, a regular positive branch,
and the named parameter path. P12--P13 retain transversality and nonlinear
nondegeneracy. A label pair alone does not prove simple eigenvalues or Hopf.
\end{itemize}

\section{Exhaustive source-range review}
Every source line from 1 through 573 is assigned below; blank separators
and environment delimiters are grouped with their substantive passage.
\textbf{PASS} concerns the stated mathematical or scope content.
\textbf{UNRESOLVED} in U01 concerns only historical numerical-run provenance.
There is no required MINOR ISSUE, MAJOR ISSUE, or FALSE mathematical finding.

{\small
\begin{longtable}{L{0.10\textwidth}L{0.17\textwidth}L{0.64\textwidth}}
\toprule P04 lines & Assessment & Finding\\\midrule\endhead
1--13 & \textbf{PASS} & Metadata, verified prerequisites, known supporting role, and scope are accurate. The declared standard ODE input is explicit.\\
14--24 & \textbf{PASS} & Canonical $n,q$, positive parameters, domain, shift convention, and absence of a coprimality hypothesis agree with current inputs.\\
25--40 & \textbf{PASS} & P03 Jacobian and all four coefficients are copied correctly. State-versus-equilibrium scope and signs of $C,G,H$ are correct.\\
41--66 & \textbf{PASS} & Complete spectrum theorem, normalized column order, both diagonalizations, and degeneracy quantifiers are correct. Local $Q$ alias is explicit (optional note T01).\\
67--85 & \textbf{PASS} & Every integer-shift phase, modular wraparound, and the separate $G$ and $-H$ orientations are correct.\\
86--104 & \textbf{PASS} & Geometric-sum normalization, square-matrix inverse argument, characteristic product, and completeness are valid.\\
105--123 & \textbf{PASS} & Both dispersion formulas pass; the imaginary $+H$ sign and signed-frequency distinction are explicit.\\
124--137 & \textbf{PASS} & Modular conjugate label and whole-vector/eigenvalue overlines are unambiguous and mathematically correct.\\
138--145 & \textbf{PASS} & Singleton counts, parity, real singleton values, and the failure of the converse are correct.\\
146--159 & \textbf{PASS} & Multiplicity statement correctly concerns numerical eigenvalues over $\CC$, with no simplicity assumption.\\
160--167 & \textbf{PASS} & Characteristic product and free-coordinate argument prove equal algebraic/geometric multiplicities and semisimplicity.\\
168--176 & \textbf{PASS} & Conjugate labels, numerical equality, real-part ties, and singleton simplicity are properly distinguished.\\
177--181 & \textbf{PASS} & Exact positive equilibrium gives $d=-4,C=2,G=H=1$ and repeated real conjugate-label values.\\
182--187 & \textbf{PASS} & Origin reduction matches verified P03; $c>0$ and distinct roots of unity imply simple values.\\
188--194 & \textbf{PASS} & Collision discussion uses the exhaustive P03 list without requiring a replacement theorem.\\
195--203 & \textbf{PASS} & $q=1$ combined matrix, spectrum, real part, and imaginary part all pass.\\
204--213 & \textbf{PASS} & $q=n-1$ phases, combined matrix, spectrum, and both dispersion reductions all pass.\\
214--223 & \textbf{PASS} & Half-offset phase $(-1)^m$, combined matrix, and zero offset imaginary contribution all pass.\\
224--228 & \textbf{PASS} & Cancellation is handled by addition, never overwriting; the general theorem remains valid.\\
229--241 & \textbf{PASS} & Zero and even singleton formulas, both $q$ parities, and the half-offset intersection agree.\\
242--259 & \textbf{PASS} & Unitary proof establishes normality without distinctness.\\
260--268 & \textbf{PASS} & Adjoint signs and powers are correct; commuting powers give a second proof. Negative states and nonsynchronous exclusions are explicit.\\
269--288 & \textbf{PASS} & Direct scalar derivative equals the unsimplified row sum at every synchronous state.\\
289--297 & \textbf{PASS} & Tangent-space differentiation and nonstationary scalar variational coefficient are correct.\\
298--309 & \textbf{PASS} & Equilibrium and positivity premises precede the optional diagonal simplification.\\
310--324 & \textbf{PASS} & Both zero-mode equilibrium expressions are equivalent; negative nonzero equilibria and origin are distinguished.\\
325--337 & \textbf{PASS} & Equilibrium scope, complete-spectrum abscissa, and exponential-stability definition are sound.\\
338--355 & \textbf{PASS} & All-label criterion, exact norm, and three linear stability cases are correctly stated.\\
356--372 & \textbf{PASS} & Exponential series, diagonal norm attainment, $t\ge0$ maximum, and converse logarithmic limit are valid.\\
373--384 & \textbf{PASS} & Real/complex norm equality and orthonormal real-pair block have the correct signs.\\
385--392 & \textbf{PASS} & Critical blocks are bounded but nondecaying; positive blocks yield instability. No Jordan growth is possible.\\
393--400 & \textbf{PASS} & Hyperbolicity and invertibility use distinct, correct all-label tests.\\
401--410 & \textbf{PASS} & Canonical origin example has a nonzero imaginary pair and no zero eigenvalue, proving strict weakness of invertibility.\\
411--414 & \textbf{PASS} & Hyperbolicity does not imply stability; nonhyperbolicity can coexist with instability.\\
415--426 & \textbf{PASS} & Nonlinear implications distinguish negative, positive, and zero abscissa and state the open-domain scope.\\
427--443 & \textbf{PASS} & Differentiable remainder, radius selection, Gronwall estimate, first-exit argument, and compact continuation all pass.\\
444--452 & \textbf{PASS} & The standard instability theorem is correctly stated and applied, including center spectrum; A04 proves this instance independently.\\
453--460 & \textbf{PASS} & Opposite cubic scalar examples establish the general limitation; they are not claimed as DM boundary examples.\\
461--472 & \textbf{PASS} & Local/global and open-domain/orthant distinctions are preserved; rejected blanket nonlinear/global claims remain rejected.\\
473--488 & \textbf{PASS} & Dominant-label ties, branch/path prerequisites, actual sign crossing, and no selection claim are properly separated.\\
489--503 & \textbf{PASS} & Least-return proof gives gcd cycles and residue classes. The full step-one coupling prevents the asserted subgraph partition from becoming a claimed full-network decomposition.\\
504--514 & \textbf{PASS} & Grid/parity/phase dependence is distinguished from coprimality independence of the proofs; the origin threshold is correctly identified as special.\\
515--529 & \textbf{PASS} & Summary describes the analytic checks actually present and does not appeal to finite computation as proof.\\
530--545 & \textbf{UNRESOLVED} & U01: historical checker execution, counts, and floating-point maxima cannot be reproduced from the supplied archive because its separate companion is absent. No theorem premise is affected.\\
546--548 & \textbf{PASS} & Explicitly limits computations to corroboration and denies that they supply an audit verdict.\\
549--573 & \textbf{PASS} & Priority checklist identifies valid audit gates and scope limits. It is reviewed as source content, not an instruction overriding the user request.\\
\bottomrule
\end{longtable}
}

\section{Issue register and closure requirements}
\textbf{Required mathematical corrections: none.}
No Fourier sign, normalization, conjugation, collision, multiplicity,
normality, stability, or hypothesis correction is required.

\textbf{U01 --- UNRESOLVED provenance, P04 \lin{530--545}.}
The old checker and per-case report are not in the supplied archive.
The exact old counts and floating-point maxima therefore remain
unreproduced. To certify those historical execution facts, obtain the
named original support package and reproduce or authenticate its records.
No P04 theorem depends on this step, and it is not a blocker for P05's
mathematical dependency gate. New exact checks must not be represented
as reproduction of that old run.

\textbf{T01 --- optional notation consistency, P04 \lin{47}.}
The canonical Fourier macro renders as $\mathsf Q$, whereas P04 introduces
the explicit local alias $Q:=\mathsf Q$. P02 has used $Q(s)$ for a scalar
polynomial. The types and the explicit local definition remove any ambiguity
in the actual P04 argument, and that polynomial is not used as a matrix
here. Retaining $\mathsf Q$ throughout would improve cross-module consistency
but is not required for correctness or closure. There is no misplaced
conjugation bar. The signed $\omega_m$ is explicitly distinguished from
future positive $\omega$.

Administrative closure may register A04 and its exact audited P04 hash
without a mathematical revision. This audit itself leaves all project
status files and P04 bytes unchanged. P05 is cleared after that closure,
while the remaining dependency gates for P08 and P11 still apply.
Neither a later proof module nor another literature survey was begun.
\end{document}
